% === PREAMBLE + SLIDES 01-02 ===
% Standalone Hetis: 10/21/2026 Research Meeting
% Source style: user's light Madrid Beamer template (same palette and components).
% Compile with pdfLaTeX. No additional font packages or shell-escape needed.
% Keep the Hetis/ figure folder next to this file.
% If an image is absent, a labeled placeholder is shown; the PDF still compiles.
% ThunderServe reviewed against the supplied arXiv:2502.09334v2 (6 Nov 2025).
% Hetis reviewed against the supplied arXiv:2509.08309v1 (10 Sep 2025).
% Independent deck. Footer numbers are 1--39 for this file only.
\documentclass[aspectratio=169,10pt]{beamer}
\usepackage[english]{babel}
\usepackage[utf8]{inputenc}
\usepackage[T1]{fontenc}
\usepackage{lmodern}
\usepackage{amsmath,amssymb}
\usepackage{booktabs,tabularx,array}
\usepackage{graphicx}
\graphicspath{{./}{figures/}}
\usepackage{xcolor}
\usepackage{tikz}
\usetikzlibrary{arrows.meta,positioning,calc,shapes.geometric}
\usepackage{hyperref}

% Light Madrid-inspired palette from the provided template.
\usetheme{Madrid}
\usecolortheme{default}
\definecolor{Ink}{HTML}{243748}
\definecolor{Muted}{HTML}{667B8B}
\definecolor{Accent}{HTML}{347E74}
\definecolor{AccentBlue}{HTML}{507DAD}
\definecolor{SoftMint}{HTML}{E6F4EE}
\definecolor{SoftBlue}{HTML}{E9F2F9}
\definecolor{SoftLilac}{HTML}{F1EEF8}
\definecolor{Paper}{HTML}{FBFCFE}
\definecolor{Divider}{HTML}{D8E4EB}

\setbeamercolor{background canvas}{bg=Paper}
\setbeamercolor{normal text}{fg=Ink}
\setbeamercolor{structure}{fg=Accent}
\setbeamercolor{frametitle}{fg=Ink,bg=SoftBlue}
\setbeamercolor{title}{fg=Ink}
\setbeamercolor{subtitle}{fg=Muted}
\setbeamercolor{block title}{fg=Ink,bg=SoftBlue}
\setbeamercolor{block body}{fg=Ink,bg=white}
\setbeamercolor{item}{fg=Accent}
\setbeamercolor{caption name}{fg=Muted}
\setbeamercolor{footline}{fg=Muted,bg=Paper}
\setbeamerfont{frametitle}{size=\Large,series=\bfseries}
\setbeamerfont{footline}{size=\tiny}
\setbeamerfont{block title}{series=\bfseries}
\setbeamertemplate{headline}{}
\setbeamertemplate{navigation symbols}{}
\setbeamertemplate{blocks}[rounded][shadow=false]
\setbeamertemplate{itemize items}[circle]
\setbeamertemplate{itemize subitem}[circle]
\setbeamertemplate{footline}{%
  \leavevmode\hbox{%
    \begin{beamercolorbox}[wd=\paperwidth,ht=.30cm,dp=.16cm,rightskip=.60cm plus 1fil]{footline}%
      \hfill\insertframenumber\,/\,\inserttotalframenumber
    \end{beamercolorbox}}}
\setbeamertemplate{frametitle}{%
  \nointerlineskip
  \begin{beamercolorbox}[wd=\paperwidth,ht=.77cm,dp=.18cm,leftskip=.64cm]{frametitle}%
    \usebeamerfont{frametitle}\insertframetitle
  \end{beamercolorbox}}

% Reusable design components: compact and easy to edit after copying.
\newcommand{\smalllabel}[1]{{\scriptsize\bfseries\color{Accent}\MakeUppercase{#1}}}
\newcommand{\hint}[1]{{\scriptsize\color{Muted}#1}}
\newcommand{\takeaway}[1]{%
  \begin{beamercolorbox}[rounded=true,sep=0.14cm,wd=\linewidth]{block title}%
    {\small\bfseries #1}
  \end{beamercolorbox}}
\newcommand{\minibadge}[2]{%
  \tikz[baseline=(node.base)]\node[fill=#1,rounded corners=3pt,
    inner xsep=6pt,inner ysep=3pt,align=center] (node) {\scriptsize\bfseries #2};}
\newcommand{\stepcard}[3]{%
  \begin{beamercolorbox}[rounded=true,sep=0.12cm,wd=\linewidth]{block body}%
    {\scriptsize\bfseries\color{Accent}#1}\par\vspace{.6mm}
    {\footnotesize\bfseries #2}\par\vspace{.6mm}
    {\scriptsize\color{Muted}#3}%
  \end{beamercolorbox}}

% Automatically use image when found, otherwise show a filename placeholder.
% This supports both main.tex's folder and figures/ (case-sensitive).
\newcommand{\figureplaceholder}[3]{%
  \begingroup\centering
  \IfFileExists{#1}{%
    \includegraphics[width=.96\linewidth,height=#2,keepaspectratio]{#1}%
  }{%
    \IfFileExists{figures/#1}{%
      \includegraphics[width=.96\linewidth,height=#2,keepaspectratio]{figures/#1}%
    }{%
      \begin{tikzpicture}
        \node[draw=AccentBlue!55,dashed,line width=.7pt,rounded corners=4pt,
              fill=SoftBlue!60,minimum height=#2,minimum width=.95\linewidth,
              text width=.85\linewidth,align=center,inner sep=2mm] {
           {\color{AccentBlue}\bfseries FIGURE PLACEHOLDER}\\[2.2mm]
           {\color{Ink}\small #3}\\[1.5mm]
           {\color{Muted}\scriptsize\ttfamily\detokenize{#1}}};
      \end{tikzpicture}%
    }%
  }\par\endgroup}

% Select layout from the original image ratio: >= 1.8 is wide.
\newsavebox{\figuremeasure}
\newcommand{\figurelayout}[3]{%
  \IfFileExists{#1}{%
    \sbox{\figuremeasure}{\includegraphics{#1}}%
    \ifdim\wd\figuremeasure<1.8\ht\figuremeasure
      #3%
    \else
      #2%
    \fi
  }{#2}}

% Wide figures: image above text. Narrow figures: text left, image right.
\newcommand{\paperfigure}[5]{%
  \figurelayout{#1}{%
    \vspace{.08cm}
    \figureplaceholder{#1}{3.20cm}{#2}
    {\centering\hint{#2}\par}
    \vspace{.06cm}
    \begin{columns}[T,onlytextwidth]
      \begin{column}{.49\textwidth}
        {\footnotesize\begin{itemize}
          \setlength{\itemsep}{.04cm}#3
        \end{itemize}}
      \end{column}
      \begin{column}{.49\textwidth}
        {\footnotesize\begin{itemize}
          \setlength{\itemsep}{.04cm}#4
        \end{itemize}}
      \end{column}
    \end{columns}
    \vspace{.10cm}\takeaway{#5}%
  }{%
    \vspace{.18cm}
    \begin{columns}[c,onlytextwidth]
      \begin{column}{.44\textwidth}
        {\footnotesize\raggedright\begin{itemize}
          \setlength{\itemsep}{.12cm}#3 #4
        \end{itemize}}
        \vspace{.15cm}\takeaway{#5}
      \end{column}
      \begin{column}{.54\textwidth}
        \figureplaceholder{#1}{5.35cm}{#2}
        \vspace{.06cm}{\centering\hint{#2}\par}
      \end{column}
    \end{columns}%
  }}

% Short discussion prompts use the same aspect-ratio-aware layout.
\newcommand{\figurefocus}[5]{%
  \paperfigure{#1}{#2}{\item #3}{\item #4}{#5}}

% Two source figures share a slide without cropping either figure.
\newcommand{\twofigurefocus}[7]{%
  \vspace{.10cm}
  \begin{columns}[c,onlytextwidth]
    \begin{column}{.49\textwidth}
      \figureplaceholder{#1}{3.78cm}{#2}
    \end{column}
    \begin{column}{.49\textwidth}
      \figureplaceholder{#3}{3.78cm}{#4}
    \end{column}
  \end{columns}
  \vspace{.18cm}
  \begin{columns}[T,onlytextwidth]
    \begin{column}{.488\textwidth}{\small\textbf{\color{Accent}01}\quad #5\par}\end{column}
    \begin{column}{.488\textwidth}{\small\textbf{\color{Accent}02}\quad #6\par}\end{column}
  \end{columns}
  \vspace{.23cm}\takeaway{#7}}

\title[Hetis]{Hetis: Fine-grained and Dynamic Parallelism}
\subtitle{LLM Serving in Heterogeneous GPU Clusters}
\author{Yu-Jui Liang}
\institute{National Taiwan University \quad | \quad Academia Sinica}
\date{October 21, 2026}
\begin{document}



% ----- SLIDE 01: Title -----
{\setbeamertemplate{footline}{}
\begin{frame}[plain]
  \vspace*{1.00cm}
  \smalllabel{Research meeting}\par
  \vspace{.61cm}
  {\raggedright\color{Ink}
    \resizebox{\textwidth}{!}{\bfseries Hetis: Fine-grained and Dynamic Parallelism}\par}
  \vspace{.39cm}
  {\color{Muted}\large LLM Serving in Heterogeneous GPU Clusters\par}
  \vspace{1.08cm}
  {\bfseries Yu-Jui Liang}\par\vspace{.05cm}
  {\small\color{Muted}Advisor: Hsiang-Yun Cheng}\par\vspace{.05cm}
  {\small\color{Muted}National Taiwan University \quad | \quad Academia Sinica \quad | \quad 21 Oct 2026\par}
  \vfill
  {\scriptsize\color{Muted}\raggedright Zizhao Mo, Jianxiong Liao, Huanle Xu, Zhi Zhou, Chengzhong Xu\par}
  {\scriptsize\color{Muted}SC 2025\par}
  \vspace{.06cm}
  {\color{Accent}\rule{\textwidth}{.7pt}}
\end{frame}}


% ----- SLIDE 02: Presentation Roadmap -----
\begin{frame}{Presentation Roadmap}
  \smalllabel{Hetis / Overview}\par
  \vspace{.22cm}
  \begin{enumerate}
    \setlength{\itemsep}{.24cm}
    \item \textbf{Background / Motivation}\par
      {\small\color{Muted}KV cache imbalance and module-specific performance gaps.}
    \item \textbf{Methodology}\par
      {\small\color{Muted}Primary parallelism, head allocation and KV cache management.}
    \item \textbf{Experimental Results}\par
      {\small\color{Muted}Latency, usable cache capacity, sensitivity and overhead.}
    \item \textbf{Discussion / Conclusion}\par
      {\small\color{Muted}Main findings, comparison with ThunderServe and research directions.}
  \end{enumerate}
  \vspace{.20cm}\hint{Source: arXiv:2509.08309 (Hetis); ThunderServe in the closing comparison.}
\end{frame}

% === HETIS SLIDES 03-05 ===
\section{Hetis}\subsection{Background / Motivation}

% ----- SLIDE 03: Coarse Splitting Wastes KV Cache Capacity -----
\begin{frame}{Coarse Splitting Wastes KV Cache Capacity}
  \smalllabel{Hetis / Background \& Motivation}\par
  \paperfigure{Hetis/Hetis_Fig01_plot.pdf}{Fig.~1; Section 2.3: Two sources of memory inefficiency}{
    \item Phase splitting stores separate model copies and confines decode KV cache to decode workers.
    \item In the FP16 7B example, a 24 GB 3090 has only about 10 GB left for KV cache.

  }{
    \item Parameter splitting balances compute, but fixes where each request stores its KV cache.
    \item One GPU can run out of cache while another still has free memory.

  }{Balancing computation alone does not guarantee usable KV cache capacity.}
\end{frame}

% ----- SLIDE: Table 1 GPU gaps and the 2.7B partition example -----
\begin{frame}{Measured GPU Gaps and a Partition Example}
  \smalllabel{Hetis / Background \& Motivation}\par

  \vspace{.08cm}
  {\footnotesize\textbf{Table 1.} OPT-2.7B, one iteration through all layers. Prefill batch 3; decode batch 25.}\par
  \vspace{.06cm}
  {\centering\small\begin{tabular}{lccc}\toprule
    Device & Memory & Time (Prefill) & Time (Decode)\\\midrule
    A100 & 80 GB & 0.06 s & 0.0097 s\\
    3090 & 24 GB & 0.147 s & 0.0143 s\\
    P100 & 12 GB & 1.47 s & 0.077 s\\\bottomrule
  \end{tabular}\par}
  \vspace{.10cm}
  {\footnotesize A100 vs.\ 3090 / P100: prefill \(2.45\times\) / \(24.5\times\); decode \(1.47\times\) / \(7.93\times\); capacity \(3.33\times\) / \(6.67\times\).}\par
  \vspace{.08cm}
  {\footnotesize\textbf{Section 2.3 example.} Split FP16 OPT-2.7B by decode compute: about 40\% of parameters and KV cache on the 3090, 60\% on the A100. The 3090 then has at most 18.6 GB for KV cache, while the A100 still has 44.1 GB unused.}\par
  \vspace{.10cm}\takeaway{The 10 GB figure on the previous slide is the FP16 7B case, not this 2.7B partition.}
  \vspace{.06cm}\hint{Source: Section 2.2, Table 1, and Section 2.3. Ratios are the paper's comparisons of these cells.}
\end{frame}

% ----- SLIDE 04: MLP and Attention Scale Differently -----
\begin{frame}{MLP and Attention Scale Differently}
  \smalllabel{Hetis / Background \& Motivation}\par
  \paperfigure{Hetis/Hetis_Fig02_plot.pdf}{Fig.~2: One Llama-70B decode layer; context length 1000}{
    \item The P100/A100 execution-time gap is much larger for MLP than for attention.
    \item Adding slow GPUs to dense operations can cost more communication than it saves in compute.

  }{
    \item Decode attention has a smaller performance gap and can use otherwise idle GPU memory.
    \item The preferred GPU partition can differ across modules within the same layer.

  }{Use selected GPUs for dense work and a wider pool for decode attention.}
\end{frame}

% ----- SLIDE 05: Primary Workers and Attention Workers -----
\begin{frame}{Primary Workers and Attention Workers}
  \smalllabel{Hetis / Methodology}\par

  {\small\setlength{\tabcolsep}{4pt}
  \begin{tabularx}{\textwidth}{@{}p{1.15cm}>{\raggedright\arraybackslash}X>{\raggedright\arraybackslash}p{3.1cm}>{\raggedright\arraybackslash}p{3.1cm}@{}}\toprule
  Phase & Operation & Primary workers & Attention workers\\\midrule
  Prefill & QKV / Output Projections, MLP & Execute & ---\\
  Prefill & Attention Core \& KV Cache Creation & Execute & ---\\\midrule
  Decode & QKV / Output Projections, MLP & Execute & ---\\
  Decode & Attention Core & Process assigned heads & Process offloaded heads\\
  Decode & KV Cache Storage & Store assigned heads & Store offloaded heads\\\bottomrule
  \end{tabularx}}\par
  \vspace{.08cm}
  {\centering \(\displaystyle\text{Attention core}(Q,K,V)=\operatorname{softmax}(QK^\top/\sqrt{d_h})V\)\par}
  \vspace{.08cm}
  {\small Here, the attention core is parameter-free \emph{after} Q, K and V are produced; the projections still require model weights.}\par
  \vspace{.12cm}\takeaway{Primary workers execute both prefill and decode, while attention workers only assist with decode attention and KV cache storage.}
  \vspace{.10cm}\hint{Source: Sections 2.4, 3.2 and 4. Attention core terminology clarified for this presentation.}
\end{frame}

% === HETIS SLIDES 06-11 ===
\subsection{Parallelization Design}

% ----- SLIDE 06: Hetis System Architecture -----
\begin{frame}{Hetis System Architecture}
  \smalllabel{Hetis / Methodology}\par
  \paperfigure{Hetis/Hetis_Fig03_plot.pdf}{Fig.~3: Parallelizer, Profiler, Dispatcher and Hauler}{
    \item Parallelizer chooses Primary GPUs and their DP / PP / TP configuration.
    \item Profiler estimates GPU computation and communication cost.

  }{
    \item Dispatcher assigns heads to Primary and Attention workers, and selectively re-dispatches existing requests.
    \item Hauler manages head-level KV blocks and migrates the required cache shards.

  }{Offline placement and profiling support online head allocation.}
\end{frame}

% ----- SLIDE 07: Hierarchical Search for Primary Parallelism -----
\begin{frame}{Hierarchical Search for Primary Parallelism}
  \smalllabel{Hetis / Methodology}\par
  \paperfigure{Hetis/Hetis_Fig04_plot.pdf}{Fig.~4; Section 4.1: Primary worker configuration search}{
    \item \textbf{1. DP grouping:} divide each GPU type evenly across serving instances; filter insufficient KV capacity.
    \item \textbf{2. Coarse PP search:} treat each GPU type as one stage; assign layers to minimize the bottleneck compute cost.

  }{
    \item \textbf{3. GPU pruning:} try low-end GPUs first. Re-evaluate layer assignments on the remaining GPUs before the 5\% test.
    \item \textbf{4. Fine-grained TP/PP search:} explore within each stage; score computation plus communication.

  }{Excluded GPUs become a shared Attention-worker pool, rather than being left unused.}
\end{frame}

% ----- SLIDE 08: Primary Selection Objective and Pruning Rule -----
\begin{frame}{Primary Selection Objective and Pruning Rule}
  \smalllabel{Hetis / Methodology}\par


  \vspace{.03cm}
  {\small\textbf{Objective (paper Eq.~1):} choose placement and DP / PP / TP configuration $\sigma$.}
  \par\vspace{.03cm}{\centering \(\displaystyle \sigma^*=\arg\min_{\sigma}C(\sigma,M,\mathcal R),\qquad C=C_{\mathrm{comp}}+C_{\mathrm{comm}}\)\par}\vspace{.03cm}
  {\footnotesize $M$: model; $\mathcal R$: workload distribution.\quad\textbf{Pipeline bottleneck (explanatory):}}
  \par\vspace{.03cm}{\centering \(\displaystyle C_p=\max_s C_{\mathrm{comp},s}\)\par}\vspace{.03cm}
  {\footnotesize $s$: coarse PP stage. Ideal scaling; computation only, no communication. This restates the text, not a separately numbered paper equation.}\par
  \vspace{.10cm}
  {\small\textbf{GPU pruning (paper rule):} re-evaluate layer assignments after removing $\kappa$.}
  \par\vspace{.03cm}{\centering \(\displaystyle \frac{C_p(\sigma-\kappa,M,\mathcal R)}{C_p(\sigma,M,\mathcal R)}\leq1+\Delta,\qquad\Delta=0.05.\)\par}\vspace{.03cm}
  {\footnotesize\textbf{Illustrative Example (assumed stage costs, arbitrary units):}
  Before: $C_p=\max(100,90,20)=100$. After removing P100 and reassigning layers: $C_p=\max(103,98)=103$. Ratio $1.03\leq1.05$: accept. Not measured results.}\par
  \vspace{.12cm}\takeaway{Reassigning layers changes the bottleneck cost and the pruning decision.}
  \vspace{.06cm}\hint{Source: Section 4.1, Eq.~(1), Figure 4 and unnumbered pruning rule. Heuristic search; no global-optimality guarantee.}
\end{frame}

% ----- SLIDE 09: Why Split Attention by Head? -----
\begin{frame}{Why Split Attention by Head?}
  \smalllabel{Hetis / Methodology}\par
  \paperfigure{Hetis/Hetis_Fig05_plot.pdf}{Fig.~5: Llama-70B, 100 Gbps; communication overhead}{
    \item Batch-wise splitting moves whole requests, which is coarse when context lengths differ.
    \item Sequence-wise splitting replicates query vectors across workers holding different token ranges.

  }{
    \item Head-wise splitting transfers only the assigned head vectors and their results.
    \item Reported communication speedups over sequence splitting: $2.68\times$ at 20\% offload; up to $3.55\times$ with four workers.

  }{Independent heads reduce communication and allow finer cache placement.}
\end{frame}

% ----- SLIDE: Eq. (2) head-wise attention -----
\begin{frame}{Head-wise Attention as Printed}
  \smalllabel{Hetis / Methodology}\par

  \vspace{.18cm}
  \[
    \mathrm{Attention}_j=\mathrm{Concat}(\mathrm{result}_{1,j},\ldots,\mathrm{result}_{n,j})
  \]
  \[
    \mathrm{result}_{i,j}
    =\operatorname{softmax}\!\left(
      q_{h_i^j(t)}\,K_{h_i^j(t)}^{\top}/\sqrt{d}
    \right)V_{h_i^j(t)}
  \]
  \vspace{.08cm}
  {\small \(h_i^j(t)\): query heads of request \(j\) placed on GPU \(i\) at time \(t\). Each GPU finishes its own heads; outputs are concatenated in head order.}\par
  \vspace{.16cm}
  \begin{itemize}\small
    \item This is Eq.~(2). The earlier attention-core line is the same softmax, without the per-GPU split.
    \item No cross-GPU softmax reduction is required inside one head.
  \end{itemize}
  \vspace{.12cm}\takeaway{Head-wise splitting changes where each head runs, not the attention definition.}
  \vspace{.06cm}\hint{Source: Section 4.2, Eq.~(2).}
\end{frame}

% ----- SLIDE 10: Dynamic Attention Placement Across Requests -----
\begin{frame}{Dynamic Attention Placement Across Requests}
  \smalllabel{Hetis / Methodology}\par
  \paperfigure{Hetis/Hetis_Fig06_plot.pdf}{Fig.~6: Three requests, six heads each; context lengths 6, 9 and 3}{
    \item Each request can use a different head distribution across the same GPU pool.
    \item KV cache follows the assigned heads; longer contexts require more cache per head.

  }{
    \item Primary and Attention workers compute their assigned heads, then gather the head outputs.
    \item Re-dispatching changes selected assignments and migrates only the cache that must move.

  }{Dynamic parallelism changes placement, not the number of model heads or the attention definition.}
\end{frame}

% ----- SLIDE 11: Head Integrity and Grouped-Query Attention -----
\begin{frame}{Head Integrity and Grouped-Query Attention}
  \smalllabel{Hetis / Methodology}\par

  \vspace{.16cm}
  \[r=\frac{H_Q}{H_{KV}},\qquad \sum_{i=1}^{N_k}x_i^j=H_Q,\qquad x_i^j/r\in\mathbb Z_{\geq0}.\]
  {\small $x_i^j$: query heads of request $j$ assigned to GPU $i$ in pipeline stage $k$.}\par
  \vspace{.18cm}
  \begin{itemize}\small
    \setlength{\itemsep}{.12cm}
    \item MHA: $r=1$; each query head has its own key/value head.
    \item GQA: each KV head is shared by a group of $r$ query heads; Hetis preserves whole groups.
    \item Illustrative example: $H_Q=8$, $H_{KV}=2$, so $r=4$; valid per-GPU allocations are 0, 4 or 8 query heads.
    \item Each assigned head uses its full context. Head-wise computation needs no cross-worker softmax reduction over sequence shards.
  \end{itemize}
  \vspace{.15cm}\takeaway{GQA reduces KV volume, but also makes the placement unit coarser than one query head.}
  \vspace{.08cm}\hint{Source: Sections 4.2, 5.2.1 and 6. Example is explanatory, not a measured model configuration.}
\end{frame}

% === HETIS SLIDES 12-18 ===
\subsection{Online Optimization and Runtime}

% ----- SLIDE 12: Profiling Decode Attention Time -----
\begin{frame}{Profiling Decode Attention Time}
  \smalllabel{Hetis / Methodology}\par
  \paperfigure{Hetis/Hetis_Fig07_plot.pdf}{Fig.~7: OPT-30B attention micro-benchmarks}{
    \item With total head count and cache size fixed, changing the number of requests has little effect in this test.
    \item Attention time grows approximately linearly with cache size.

  }{
    \item At a fixed cache size, more query heads also increase execution time.
    \item These observations motivate a per-GPU affine model using head count and cache amount.

  }{The model uses the work assigned to a GPU, not request count alone.}
\end{frame}

% ----- SLIDE 13: Computation and Communication Cost Models -----
\begin{frame}{Computation and Communication Cost Models}
  \smalllabel{Hetis / Methodology}\par

  \vspace{.13cm}
  \[\tau_i=a_i h_i+b_i g_i+c_i,\qquad \rho_i=\gamma_i d_i+\beta_i,\qquad d_i=(2+2/r)h_i.\]
  \begin{itemize}\small
    \setlength{\itemsep}{.13cm}
    \item $h_i$: total query heads across requests; $g_i$: associated KV cache amount.
    \item $a_i,b_i,c_i$: fitted compute coefficients; $\gamma_i,\beta_i$: fitted transfer coefficients.
    \item Per decode step, remote heads need Q, new K/V and returned attention outputs: $2h_i$ query/output vectors plus $2h_i/r$ KV vectors.
    \item The expression for $d_i$ counts vectors; head width and bytes per element must be accounted for when using physical byte units.
  \end{itemize}
  \vspace{.15cm}\takeaway{Attention-worker cost includes network transfer; retaining work on a Primary GPU avoids that extra transfer.}
  \vspace{.08cm}\hint{Source: Section 5.1, Eqs.~(3)--(4). Vector-count interpretation makes the units explicit.}
\end{frame}

% ----- SLIDE 14: Online Dispatch for Newly Arriving Requests -----
\begin{frame}{Online Dispatch for Newly Arriving Requests}

\smalllabel{Hetis / Methodology}\par
\vspace{.10cm}
{\small Section 5.2.2 keeps dispatched placements fixed and assigns only new requests $\mathrm J(t)$.}\par
\vspace{.12cm}
{\small\textbf{Eq. (8): load update after solving the dispatch problem}}
\[
\begin{aligned}
h_i(t+1)&=h_i(t)+\sum_{j=1}^{\mathrm J(t)}x_i^j(t),\\
g_i(t+1)&=g_i(t)+\frac{2}{r}\sum_{j=1}^{\mathrm J(t)}x_i^j(t)\,l_j(t).
\end{aligned}
\]
\begin{itemize}\small
\item $x_i^j(t)$: query heads of request $j$ assigned to GPU $i$ in stage $k$; $N_k$: workers in that stage.
\item $h_i(t)$: existing query-head count; $g_i(t)$: existing cache amount; $l_j(t)$: request context length.
\item $r$: query/KV-head ratio; $2/r$ accounts for K and V shared by each group.
\end{itemize}
\vspace{.08cm}\takeaway{Eq. (8) updates dispatch load; it does not specify all cache-growth or departure updates.}
\vspace{.06cm}\hint{Source: Section 5.2.2, Eq. (8). Printed upper limit J(t) is retained; the prose defines J(t) as a set.}
\end{frame}

% ----- SLIDE: Printed Eqs. (5)--(7) -----
\begin{frame}{Printed Dispatch Constraints}

\smalllabel{Hetis / Methodology}\par
\vspace{.06cm}
{\small\setlength{\abovedisplayskip}{3pt}\setlength{\belowdisplayskip}{3pt}\setlength{\abovedisplayshortskip}{3pt}\setlength{\belowdisplayshortskip}{3pt}
\[
\begin{aligned}
&\sum\nolimits_{i=1}^{N_k}x_i^j(t)\in\{0,\mathrm H\},\quad\forall j\in\mathbb S(t),\ \forall t.&&\text{(5)}\\
&\sum\nolimits_{j\in\mathbb S(t)}x_i^j(t)l_j(t)\leq\frac{r\mathrm M_i(t)}{2},\quad\forall 1\leq i\leq N_k,\ \forall t.&&\text{(6)}
\end{aligned}
\]}
{\footnotesize\textbf{Eq. (7): only newly arriving requests}}
{\small\setlength{\abovedisplayskip}{3pt}\setlength{\belowdisplayskip}{3pt}\setlength{\abovedisplayshortskip}{3pt}\setlength{\belowdisplayshortskip}{3pt}
\[
\begin{aligned}
&\min\max_i f_i(\vec x_i(t)), &&\text{(7a)}\\
&g_i(t)+\sum\nolimits_{j=1}^{\mathrm J(t)}x_i^j(t)l_j(t)\leq\frac{r\mathrm M_i(t)}{2},\quad\forall i,t, &&\text{(7b)}\\
&\sum\nolimits_{i=1}^{N_k}x_i^j(t)=\mathrm H,\quad\forall j\in\mathrm J(t). &&\text{(7c)}
\end{aligned}
\]}
{\footnotesize $\mathbb S(t)$: in-service requests; $\mathrm J(t)$: new requests; $\mathrm H$: query heads; $\mathrm M_i(t)$: available KV capacity.}\par
\vspace{.06cm}\takeaway{Keep printed (7b) unchanged: its cache scaling differs from Eq. (8).}
\vspace{.04cm}\hint{Source: Section 5.2, Eqs. (5)--(7). Group integrality appears in the prose, not in printed (7).}
\end{frame}

% ----- SLIDE: Closed forms of f_i -----
\begin{frame}{Closed Forms of Attention Time}

\smalllabel{Hetis / Methodology}\par
\vspace{.10cm}
{\footnotesize\textbf{Primary worker: computation only (unnumbered, after Eq. (7))}}
{\small\setlength{\abovedisplayskip}{3pt}\setlength{\belowdisplayskip}{3pt}\setlength{\abovedisplayshortskip}{3pt}\setlength{\belowdisplayshortskip}{3pt}
\[
\begin{aligned}
f_i(\vec x_i(t))={}&a_i\Big(h_i(t)+\sum\nolimits_{j=1}^{\mathrm J(t)}x_i^j(t)\Big)\\
&+b_i\Big(g_i(t)+\frac{2}{r}\sum\nolimits_{j=1}^{\mathrm J(t)}l_j(t)x_i^j(t)\Big)+c_i.
\end{aligned}
\]}
{\footnotesize\textbf{Attention worker: computation plus transfer (same paragraph)}}
{\small\setlength{\abovedisplayskip}{3pt}\setlength{\belowdisplayskip}{3pt}\setlength{\abovedisplayshortskip}{3pt}\setlength{\belowdisplayshortskip}{3pt}
\[
\begin{aligned}
f_i(\vec x_i(t))={}&\Big(a_i+\Big(2+\frac{2}{r}\Big)\gamma_i\Big)
\Big(h_i(t)+\sum\nolimits_{j=1}^{\mathrm J(t)}x_i^j(t)\Big)\\
&+b_i\Big(g_i(t)+\frac{2}{r}\sum\nolimits_{j=1}^{\mathrm J(t)}l_j(t)x_i^j(t)\Big)+c_i+\beta_i.
\end{aligned}
\]}
{\footnotesize $a_i,b_i,c_i$: compute coefficients; $\gamma_i,\beta_i$: transfer coefficients. The transfer term uses existing plus newly assigned heads; it sends Q, new K/V and outputs, not the entire historical cache each step.}\par
\vspace{.08cm}\takeaway{Both expressions use the same $2/r$ cache increment as Eq. (8).}
\vspace{.04cm}\hint{Source: Section 5.2.2, unnumbered formulas after Eq. (7). Time indices and printed sum limits restored.}
\end{frame}

% ----- SLIDE 15: Dispatch Formulation: What the Paper Specifies -----
\begin{frame}{Dispatch Formulation: What the Paper Specifies}

\smalllabel{Hetis / Methodology}\par
\vspace{.10cm}
{\small\textbf{Derived: epigraph form of the minimax objective only}}
\[
\min_{x,z}z\quad\text{s.t.}\quad f_i(\vec x_i(t))\leq z\quad\forall i.
\]
\begin{itemize}\small
\setlength{\itemsep}{.09cm}
\item Retain the allocation and capacity constraints separately. This objective rewrite does not repair printed (7b).
\item Section 5.2.1 requires $x_i^j(t)/r\in\mathbb N$. The paper does not define whether $0\in\mathbb N$; Eq. (5) permits zero total allocation.
\item Continuous variables give an LP relaxation; retaining group integrality gives a discrete problem. Section 5.2.2 claims LP reformulation but does not explain the integrality treatment.
\item (7b) adds $g_i$ to query-head-token counts; Eq. (8) and $f_i$ multiply new counts by $2/r$. Cache units and normalization are not fully specified.
\end{itemize}
\vspace{.08cm}\takeaway{No claimed equivalence to a repaired capacity rule; no assumed rounding or solver algorithm.}
\vspace{.06cm}\hint{Source: Section 5.2. Presenter analysis: J(t) is used as both a set and a sum limit; physical byte conversion is unspecified.}
\end{frame}

% ----- SLIDE 16: Re-dispatching to Balance Computation -----
\begin{frame}{Re-dispatching to Balance Computation}
  \smalllabel{Hetis / Methodology}\par

  \vspace{.04cm}
  {\footnotesize\textbf{Ideal attention time over existing requests $\mathbb S(t)$:}}
  \[
    f^*=\min\max_i f_i(\vec x_i(t))
    \quad\text{s.t.}\quad
    \sum_{i=1}^{N_k}\sum_{j\in\mathbb S(t)}x_i^j(t)\,l_j(t)
    \le\sum_{i=1}^{N_k}\frac{r\cdot\mathrm M_i}{2}.
  \]
  \vspace{-.18cm}
  \begin{itemize}\small
    \setlength{\itemsep}{.05cm}
    \item If the deviation from the current attention time exceeds $\Theta$ (default 50\%), pick the slowest GPU and its heaviest request.
    \item Re-dispatch that request with Eq.~(7). Keep reusable KV shards; transfer only changed assignments.
    \item The displayed constraint is one cluster-wide sum. The paper does not define how ``deviation'' is normalized.
  \end{itemize}
  \vspace{.06cm}\takeaway{Growing contexts can create stragglers even if initial head placement was balanced.}
  \vspace{.04cm}\hint{Source: Section 5.3.1. Head-integrity constraints are not restated in this $f^*$ display.}
\end{frame}

% ----- SLIDE 17: Re-dispatching to Relieve KV Cache Exhaustion -----
\begin{frame}{Re-dispatching to Relieve KV Cache Exhaustion}
  \smalllabel{Hetis / Methodology}\par

  \vspace{.22cm}
  \begin{itemize}\small
    \setlength{\itemsep}{.20cm}
    \item A globally selected LIFO / LRU victim may hold no cache on the exhausted GPU.
    \item Restrict victim selection to requests occupying the overloaded device; the paper adapts LIFO to that device.
    \item Try redistribution when \(\sum_i g_i(t)<\sum_i r\cdot\mathrm M_i/2\); otherwise the cluster has no aggregate free cache under the paper's units.
    \item Repeat subject to remaining cache capacity and the computation-imbalance threshold; reuse heads whose placement is unchanged.
  \end{itemize}
  \vspace{.22cm}\takeaway{Target the actual memory bottleneck and migrate only the necessary KV shards.}
  \vspace{.10cm}\hint{Source: Section 5.3.2. Aggregate free capacity alone does not guarantee a feasible discrete allocation.}
\end{frame}

% ----- SLIDE 18: Head-Level KV Management and Migration -----
\begin{frame}{Head-Level KV Management and Migration}
  \smalllabel{Hetis / Methodology}\par

  \vspace{.20cm}
  \begin{columns}[T,onlytextwidth]
    \begin{column}{.49\textwidth}
      \stepcard{CACHE LAYOUT}{Split blocks along the head dimension}{Extend vLLM's fixed-sized blocks; index by sequence ID, token position and head ID.}
      \vspace{.17cm}
      \stepcard{ATTENTION KERNEL}{Reuse PagedAttention}{Compute local heads independently; accelerate block indexing with CPU cores.}
    \end{column}
    \begin{column}{.49\textwidth}
      \stepcard{COMMUNICATION}{NCCL point-to-point groups}{Transfer head-specific inputs and results between Primary and Attention workers.}
      \vspace{.17cm}
      \stepcard{LIVE MIGRATION}{Low-priority CUDA streams}{Reduce interference with ongoing inference communication; bandwidth cost still exists.}
    \end{column}
  \end{columns}
  \vspace{.22cm}\takeaway{Fine-grained placement requires matching cache indexing, kernels and transfer support.}
  \vspace{.10cm}\hint{Source: Section 6. Low-priority migration is not a guarantee of zero overhead.}
\end{frame}

% === HETIS SLIDES 19-30 ===
\subsection{Experimental Results}

% ----- SLIDE 19: Hetis Experimental Setup -----
\begin{frame}{Hetis Experimental Setup}
  \smalllabel{Hetis / Experimental Results}\par

  \vspace{.20cm}
  \begin{columns}[T,onlytextwidth]
    \begin{column}{.49\textwidth}
      \stepcard{CLUSTER}{12 GPUs across four hosts}{One host with 4 A100-80GB; two hosts with 2 RTX 3090 each; one host with 4 P100.}
      \vspace{.17cm}
      \stepcard{NETWORK}{100 Gbps LAN; intra-host PCIe}{A local heterogeneous cluster, rather than a cloud-rental cost comparison.}
    \end{column}
    \begin{column}{.49\textwidth}
      \stepcard{MODELS}{MHA and GQA}{Llama-13B and OPT-30B use MHA; Llama-70B uses GQA in the evaluation.}
      \vspace{.17cm}
      \stepcard{WORKLOADS}{SG / HE / LB}{ShareGPT: chatbot; HumanEval: code completion; LongBench: long-article summarization.}
    \end{column}
  \end{columns}
  \vspace{.20cm}\takeaway{Compare Hetis with Splitwise and HexGen under the paper's specified deployments.}
  \vspace{.10cm}\hint{Source: Section 7.1. Table 1 lists 3090 with 24 GB and P100 with 12 GB.}
\end{frame}

% ----- SLIDE 20: Baselines and Evaluation Metrics -----
\begin{frame}{Baselines and Evaluation Metrics}
  \smalllabel{Hetis / Experimental Results}\par

  \vspace{.18cm}
  \begin{itemize}\small
    \setlength{\itemsep}{.14cm}
    \item \textbf{Llama-70B deployment:} Hetis selects A100 / 3090 as Primary workers and P100 as Attention workers.
    \item \textbf{HexGen:} asymmetric parameter splitting with PP across GPU types and TP within stages.
    \item \textbf{Splitwise:} authors configure A100 TP=4 for prefill; decode uses two PP stages with TP=2 on 3090 and P100.
    \item \textbf{Figures 8--10:} normalized E2E latency in s/token vs. offered request rate. A later sharp rise indicates greater load tolerance.
    \item \textbf{Figures 12--13:} P95 latency normalized to Hetis, which is 1.0. These are ratios, unlike the s/token axes above.
  \end{itemize}
  \vspace{.16cm}\takeaway{Keep deployment assumptions and normalization conventions explicit when interpreting speedups.}
  \vspace{.09cm}\hint{Source: Section 7.2. No fixed SLO goodput threshold or exact E2E normalization formula is specified.}
\end{frame}

% ----- SLIDE 21: Llama-13B: Latency Under Increasing Load -----
\begin{frame}{Llama-13B: Latency Under Increasing Load}
  \smalllabel{Hetis / Experimental Results}\par
  \paperfigure{Hetis/Hetis_Fig08_plot.pdf}{Fig.~8: SG / HE / LB; normalized E2E latency in s/token}{
    \item Compare methods within each dataset at the same offered request rate.
    \item Hetis delays the sharp latency increase across all three workload curves.

  }{
    \item The saturation points differ by dataset; each panel has its own axis scale.
    \item The advantage is most visible as load increases and compute or cache capacity becomes limiting.

  }{The figure supports higher sustainable load, not a universal per-request latency ratio.}
\end{frame}

% ----- SLIDE 22: OPT-30B: Latency Under Increasing Load -----
\begin{frame}{OPT-30B: Latency Under Increasing Load}
  \smalllabel{Hetis / Experimental Results}\par
  \paperfigure{Hetis/Hetis_Fig09_plot.pdf}{Fig.~9: OPT-30B (MHA), three workload families}{
    \item Hetis also delays the rapid latency increase for the larger MHA model.
    \item With larger model weights, placement determines how much capacity remains for KV cache.

  }{
    \item Head allocation can use memory that static compute-based partitioning leaves underutilized.
    \item Compare curves at matched load; do not compare raw heights across different dataset panels.

  }{The benefit persists for OPT-30B in the evaluated heterogeneous cluster.}
\end{frame}

% ----- SLIDE 23: Llama-70B: GQA End-to-End Results -----
\begin{frame}{Llama-70B: GQA End-to-End Results}
  \smalllabel{Hetis / Experimental Results}\par
  \paperfigure{Hetis/Hetis_Fig10_plot.pdf}{Fig.~10: Llama-70B; Section 7.2 summarizes Figures 8--10}{
    \item Hetis retains the advantage with a GQA model, whose KV cache per query head is smaller.
    \item A100 / 3090 handle dense work; P100 contributes decode attention and cache capacity.

  }{
    \item Across the end-to-end experiments, reported throughput gains are up to $2.25\times$ vs. Splitwise and $1.33\times$ vs. HexGen.
    \item These maxima summarize the experiments; they are not guaranteed for every panel or workload.

  }{Module placement and dynamic attention sharing remain useful with GQA.}
\end{frame}

% ----- SLIDE 24: Maximum Usable KV Cache Space -----
\begin{frame}{Maximum Usable KV Cache Space}
  \smalllabel{Hetis / Experimental Results}\par
  \paperfigure{Hetis/Hetis_Fig11_plot.pdf}{Fig.~11: Maximum available KV cache during inference}{
    \item Hetis reports up to $1.87\times$ more usable KV cache capacity than the baselines.
    \item It avoids storing full model copies on attention-only workers.

  }{
    \item Dynamic allocation reduces the idle cache caused by fixed compute-based partitions.
    \item The gain is usable cache capacity, not additional physical GPU memory or a guaranteed equal throughput gain.

  }{More usable cache allows the system to host more concurrent request state.}
\end{frame}

% ----- SLIDE 25: P95 TTFT and TPOT on Llama-70B -----
\begin{frame}{P95 TTFT and TPOT on Llama-70B}
  \smalllabel{Hetis / Experimental Results}\par
  \paperfigure{Hetis/Hetis_Fig12_plot.pdf}{Fig.~12: SG 1.5 req/s; HE 6 req/s; LB 0.8 req/s}{
    \item All baselines are below saturation at these selected request rates.
    \item P95 TTFT ratio: up to $1.22\times$ vs. HexGen and $1.47\times$ vs. Splitwise.

  }{
    \item P95 TPOT ratio: up to $1.39\times$ over the baselines.
    \item Primary selection helps prefill; avoiding slow dense execution and sharing attention help decode.

  }{The tail-latency comparison tests unsaturated operation, beyond maximum load capacity.}
\end{frame}

% ----- SLIDE 26: MLP and Attention Contributions to Decode Latency -----
\begin{frame}{MLP and Attention Contributions to Decode Latency}
  \smalllabel{Hetis / Experimental Results}\par
  \paperfigure{Hetis/Hetis_Fig13_plot.pdf}{Fig.~13: Llama-70B; the same request rates as Fig.~12}{
    \item MLP P95 latency improves by up to $1.29\times$; attention by up to $1.49\times$.
    \item Dense and attention modules use different placement decisions.

  }{
    \item For this experiment, module latency is the maximum stage time multiplied by the number of PP stages.
    \item The $1.49\times$ is a module-level result here; it should not be relabeled as a universal E2E speedup.

  }{Selective dense parallelism and dynamic attention sharing contribute separately.}
\end{frame}

% ----- SLIDE 27: Dynamic Head and Cache Allocation -----
\begin{frame}{Dynamic Head and Cache Allocation}
  \smalllabel{Hetis / Experimental Results}\par
  \paperfigure{Hetis/Hetis_Fig14_plot.pdf}{Fig.~14: Llama-13B; one A100 Primary + two 3090 Attention GPUs}{
    \item ShareGPT arrival rate changes over time: $5\rightarrow0\rightarrow2.5\rightarrow0$ req/s.
    \item A100 processes more heads, while both GPU types can reach high cache utilization.

  }{
    \item 3090 participation starts later under light load because remote execution adds communication cost.
    \item The head counts and cache usage change with both arrivals and completion of ongoing requests.

  }{Hetis activates remote attention when its compute and memory benefits justify communication.}
\end{frame}

% ----- SLIDE 28: Re-dispatching Benefits and Cache Operation Costs -----
\begin{frame}{Re-dispatching Benefits and Cache Operation Costs}
  \smalllabel{Hetis / Experimental Results}\par
  \paperfigure{Hetis/Hetis_Fig15_plot.pdf}{Fig.~15: Re-dispatching vs. LIFO; head management vs. vLLM}{
    \item ShareGPT at 5 req/s: re-dispatching improves output latency relative to LIFO.
    \item P95 ratio is about $1.14\times$; mean ratio is inconsistent: text $1.06\times$, plotted bars $1.16\times$.

  }{
    \item Head-wise cache storage-operation overhead rises by 13\%; fetching time falls by 26\%.
    \item The 13\% is not 13\% more KV bytes: the paper states equal KV memory consumption.

  }{Re-dispatching has benefits and implementation costs; preserve the metric and comparison.}
\end{frame}

% ----- SLIDE 29: Sensitivity to Threshold and Profiling Error -----
\begin{frame}{Sensitivity to Threshold and Profiling Error}
  \smalllabel{Hetis / Experimental Results}\par
  \paperfigure{Hetis/Hetis_Fig16_plot.pdf}{Fig.~16: Re-dispatch threshold and fitted-model perturbations}{
    \item Default $\Theta=0.5$ balances migration frequency against computation imbalance.
    \item Smaller thresholds can migrate too often; larger thresholds can leave stragglers uncorrected.

  }{
    \item Perturbing fitted coefficients by up to $\pm20\%$ raises per-token latency by at most 6.9\% in the reported tests.
    \item This evaluates model-coefficient error, not all network failures or workload shifts.

  }{The tested sensitivity is moderate; it is not a general robustness guarantee.}
\end{frame}

% ----- SLIDE 30: Planning and Profiling Overhead -----
\begin{frame}{Planning and Profiling Overhead}
  \smalllabel{Hetis / Experimental Results}\par

  \vspace{.22cm}
  {\small\begin{tabularx}{\textwidth}{@{}p{4.0cm}Xp{3.0cm}@{}}\toprule
  Component & Setting & Reported result\\\midrule
  Primary placement search & Local heterogeneous cluster & 4 seconds\\
  Larger placement search & Simulation: 5 types $\times$ 32 GPUs & 15 seconds\\
  Profiling samples & 8 head counts $\times$ 8 cache amounts & Typically $\leq100$ ms per configuration\\
  Transfer model accuracy & Profiled settings & 92.4--96.1\%\\\bottomrule
  \end{tabularx}}\par
  \vspace{.23cm}
  \begin{itemize}\small
    \item Placement search runs before initial deployment; the 160-GPU result is a simulation.
    \item Computation-model accuracy is reported as up to 93.8\%, without a precise accuracy formula.
    \item These timings do not establish the latency distribution of the online dispatch solver.
  \end{itemize}
  \vspace{.16cm}\takeaway{Separate one-time planning cost, per-sample profiling cost and online scheduling overhead.}
  \vspace{.09cm}\hint{Source: Section 7.4.}
\end{frame}

% === HETIS SLIDES 31-32 ===
\subsection{Discussion and Conclusion}

% ----- SLIDE 31: Hetis: Scope and Open Questions -----
\begin{frame}{Hetis: Scope and Open Questions}
  \smalllabel{Hetis / Discussion}\par

  \vspace{.20cm}
  \begin{itemize}\small
    \setlength{\itemsep}{.18cm}
    \item \textbf{Granularity:} GQA groups restrict how finely heads and memory can be divided; integer handling is underspecified.
    \item \textbf{Network:} repeated per-layer input/result transfers and occasional KV migration may behave differently on slower or congested links.
    \item \textbf{Evidence:} end-to-end tests use a local 12-GPU cluster; the larger-scale result measures simulated placement search.
    \item \textbf{Reproduction:} memory-unit conventions, E2E normalization and online solver latency need more detail.
  \end{itemize}
  \vspace{.20cm}\takeaway{Research question: how should head placement change when memory capacity grows faster than memory bandwidth?}
  \vspace{.10cm}\hint{Presenter analysis based on Sections 4--7; the research question is not evaluated in Hetis.}
\end{frame}

% ----- SLIDE 32: Hetis: Main Contributions and Findings -----
\begin{frame}{Hetis: Main Contributions and Findings}
  \smalllabel{Hetis / Conclusion}\par

  \vspace{.20cm}
  \begin{itemize}\small
    \setlength{\itemsep}{.20cm}
    \item \textbf{Primary parallelism:} select GPUs and DP / PP / TP configurations for dense work, including execution through both inference phases.
    \item \textbf{Dynamic attention:} assign whole head groups per request across Primary and Attention workers using compute, network and cache models.
    \item \textbf{Runtime support:} selectively re-dispatch growing requests and migrate only necessary head-level KV shards.
    \item \textbf{Reported results:} throughput up to $2.25\times$ vs. Splitwise and $1.33\times$ vs. HexGen; usable KV capacity up to $1.87\times$.
  \end{itemize}
  \vspace{.22cm}\takeaway{Different modules need different GPU partitions, and attention placement should adapt online.}
  \vspace{.10cm}\hint{Source: Sections 3--7 and 9. Ratios are maxima in the evaluated settings.}
\end{frame}

% === CROSS-PAPER DISCUSSION SLIDES 33-35 ===
\section{Cross-Paper Discussion}

% ----- SLIDE 33: ThunderServe vs. Hetis: Design Choices -----
\begin{frame}{ThunderServe vs. Hetis: Design Choices}
  \vspace{.06cm}\smalllabel{Cross-paper Comparison}\par

  \vspace{.26cm}
  {\footnotesize
  \renewcommand{\arraystretch}{1.45}
  \begin{tabularx}{\textwidth}{@{}p{2.25cm}XX@{}}
    \toprule
    \textbf{Dimension} & \textbf{ThunderServe} & \textbf{Hetis} \\
    \midrule
    Granularity & Prefill / decode phases & Dense modules / attention heads \\
    Workers & Phase-specific model replicas & Primary and Attention workers \\
    Search & Tabu search + lower-level selection & Primary search + head dispatch \\
    Runtime changes & Phase roles and routing & Head allocation / redispatch \\
    Data movement & Inter-replica KV transfer & Per-layer Q/K/V/results; KV migration \\
    Metrics & SLO attainment, throughput & E2E s/token, P95, usable KV, request rate \\
    \bottomrule
  \end{tabularx}}
  \vspace{.32cm}
  \takeaway{ThunderServe and Hetis are not compared head-to-head in these papers.}

\vspace{.08cm}\hint{Sources: ThunderServe Sections 3--5; Hetis Sections 4--7.}
\end{frame}

% ----- SLIDE 34: Implications for GPU + HBM / HBF Co-design -----
\begin{frame}{Implications for GPU + HBM / HBF Co-design}
  \vspace{.06cm}\smalllabel{Cross-paper / Research Directions}\par

  \vspace{.25cm}
  \begin{columns}[T,onlytextwidth]
    \begin{column}{.32\textwidth}
      \stepcard{QUESTION 1}{Memory-aware placement}{How do HBM bandwidth and capacity alter phase or head assignment?}
    \end{column}
    \begin{column}{.32\textwidth}
      \stepcard{QUESTION 2}{Tiered KV caches}{Could a slower, larger tier improve capacity without hurting TPOT?}
    \end{column}
    \begin{column}{.32\textwidth}
      \stepcard{QUESTION 3}{Serving simulation}{What hardware-model extensions are needed in LLMServingSim?}
    \end{column}
  \end{columns}
  \vspace{.39cm}
  \takeaway{These are follow-up research hypotheses, not evaluated HBF results in either paper.}
  \vspace{.16cm}\hint{Discussion focus: scheduling policy, cost-model fidelity, and memory hierarchy.}

\end{frame}

% ----- SLIDE 35: Q&A -----
\begin{frame}[plain]
  \vfill
  \begin{center}
    {\fontsize{34}{42}\selectfont\bfseries\color{Ink}Questions \& Discussion}\par
    \vspace{.36cm}
    {\Large\color{Accent}GPU heterogeneity \quad / \quad KV cache \quad / \quad HBM--HBF}\par
    \vspace{.50cm}
    {\small\color{Muted}ThunderServe (phase level) \quad vs. \quad Hetis (head level)}
  \end{center}
  \vfill
\end{frame}

\end{document}
